\chapter{Array-Based Lists}\label{ch:arrays}

\section{Arrays and their shortcomings}

An array stores $m$ items of one type in $m$ consecutive memory cells, so the address of
cell $i$ is computed in one step: \cc{a[i]} costs $\Theta(1)$. The \emph{primitive}
C/C++ array has well-known shortcomings, which motivate every container in this chapter
(lecture slides, \emph{why the primitive array is problematic}):
\begin{itemize}
\item its size is not stored --- a function receiving \cc{int a[]} cannot ask how long it is;
\item bounds (subscripts) are not checked: \cc{a[m]} silently reads or corrupts a
  neighbouring cell;
\item a dynamically sized array needs manual heap allocation;
\item inserting or deleting in the middle requires manual shifting;
\item an array parameter is passed by address, so the callee changes the caller's array.
\end{itemize}
The remedy is to \emph{wrap} the array together with its bookkeeping, so that the pair
behaves as a list. The slides do this twice: the \emph{array container} keeps only the
array; the \emph{array list} keeps the array \emph{and the size}.

\section{The array list}

\begin{definition}[Array list]
An \emph{array list} keeps a fixed-capacity array \cc{items[MAX]} and an integer
\cc{size}. The list is \cc{items[0..size-1]}; invariant~(I) of
Example~\ref{ex:alinv} holds.
\end{definition}

The sample program \cc{arraylist.c} defines the structure and the operations of
Table~\ref{tab:alops}; Listing~\ref{lst:arraylist} is its completed form. The slides
show the same list as a C++ class \cc{ArrayList} with \cc{push\_back}, \cc{insert(pos,x)}
and \cc{erase(pos)}, both 0-based (Tracing~\ref{trace:4}).

\begin{table}[ht]\centering\small
\setlength{\tabcolsep}{4pt}\begin{tabular}{@{}llll@{}}\toprule
Operation & Effect & Positions & Cost\\\midrule
\cc{init(l)} & \cc{size=0} & --- & $\Theta(1)$\\
\cc{append(l,x)} / \cc{push\_back} & store at \cc{items[size]} & --- & $\Theta(1)$\\
\cc{insertFront(l,x)} & shift all right, store at \cc{items[0]} & --- & $\Theta(n)$\\
\cc{insert(pos,x)} & shift \cc{items[pos..size-1]} right & 0-based (slides) & $\Theta(n-\mathrm{pos})$\\
\cc{removeLast(l)} & \cc{size--} & --- & $\Theta(1)$\\
\cc{removeFirst(l)} & shift all left, \cc{size--} & --- & $\Theta(n)$\\
\cc{erase(pos)} & shift \cc{items[pos+1..]} left & 0-based (slides) & $\Theta(n-\mathrm{pos})$\\
\cc{itemAt(l,p)} / \cc{setItem(l,x,p)} & read / write & \textbf{1-based} (\cc{arraylist.c}) & $\Theta(1)$\\
\cc{isFull(l)} & \cc{size==MAX} & --- & $\Theta(1)$\\\bottomrule
\end{tabular}
\caption{Operations of the array list. Check the statement for the position convention.}
\label{tab:alops}
\end{table}

\listingfile[label=lst:arraylist]{listings/arraylist.c}{The array list of the sample program \cc{arraylist.c}, completed.}

\begin{figure}[ht]\centering
\begin{tikzpicture}[x=7.5mm]
\node[left,font=\small] at (-0.6,0) {before:};
\foreach \v [count=\i from 0] in {16,4,8,15,{}} \node[cell] (b\i) at (\i,0) {\v};
\foreach \i in {0,...,4} \node[idx,above=0pt of b\i] {\i};
\node[left,font=\small] at (-0.6,-1.4) {after:};
\foreach \v [count=\i from 0] in {23,16,4,8,15} \node[cell] (a\i) at (\i,-1.4) {\v};
\node[cell,sh] at (0,-1.4) {23};
\foreach \i [evaluate=\i as \j using int(\i+1)] in {0,...,3} \draw[ptr] (b\i.south) to[out=-60,in=120] (a\j.north);
\node[right,font=\scriptsize,align=left] at (5.2,-0.7) {\texttt{insertFront(l,23)}:\\ every item moves one cell right,\\ \emph{last one first}; $n$ moves.};
\end{tikzpicture}
\caption{Insertion at the front shifts the whole list.}
\end{figure}

\begin{theorem}[Cost of positional insertion and deletion]\label{thm:shift}
Inserting at (0-based) position $p$ into an array list of size $n$ performs exactly $n-p$
element moves; deleting the item at position $p$ performs exactly $n-1-p$ moves. Hence
both are $\Theta(n)$ in the worst case ($p=0$) and $\Theta(1)$ at the rear ($p=n$ or
$p=n-1$).
\end{theorem}
\begin{proof}
Insertion moves \cc{items[p..n-1]}, which are $n-p$ cells, each once. Deletion moves
\cc{items[p+1..n-1]}, which are $n-1-p$ cells.
\end{proof}

\begin{note}{Direction of the shift}
Shift to the \emph{right} (insertion) from the \emph{end} towards the position; shift to
the \emph{left} (deletion) from the position towards the end. Doing it the other way
overwrites items before they are copied and turns the whole tail into copies of one value
--- a classical tracing question. The slides show the intermediate states
\cc{1 3 3 4 \dots}, \cc{1 3 4 4 \dots}: during the shift one value is always duplicated,
and the duplicate at the end (``two items equal to 10'') is harmless because it lies
beyond \cc{size}.
\end{note}

\section{Dynamic arrays and the vector}\label{sec:vector}

A fixed \cc{MAX} fails when the list outgrows it. A \emph{dynamic array} (the
\emph{ArrayStack} of the notes, the \cc{std::vector} of C++) keeps a backing array of
length $\cc{cap}\ge n$ and replaces it when necessary.
\begin{itemize}
\item \textbf{Grow.} If an insertion finds $n+1>\cc{cap}$, allocate a new array of length
  $\max(1,2n)$, copy the $n$ items, release the old array.
\item \textbf{Shrink.} After a removal, if $\cc{cap}\ge3n$, resize to $\max(1,2n)$ so
  that memory is not wasted.
\end{itemize}

\listingfile[language=C++]{listings/intvector.cpp}{A growable array of integers (tested against an array model on random operations).}

Twenty \cc{push\_back} calls followed by twenty \cc{pop\_back} calls print the pairs
\emph{size/capacity}:
\outputfile{vector_demo}
The capacity doubles exactly when the array is full, and shrinks only when it is at
least three times the size.

\begin{theorem}[Amortised cost of \cc{push\_back}]\label{thm:amortised}
Starting from an empty dynamic array that doubles when full, any sequence of $m$ calls to
\cc{push\_back} costs $O(m)$ in total, i.e.\ $O(1)$ amortised per call.
\end{theorem}
\begin{proof}
A call that does not resize costs 1. Resizes occur when the size is $1,2,4,\dots,2^k$
with $2^k<m$, and the resize at size $2^j$ copies $2^j$ items. The total copying cost is
$\sum_{j=0}^{k}2^j=2^{k+1}-1<2m$. Adding the $m$ unit costs gives a total below $3m$.
\end{proof}

\begin{remark}[Why shrink at $3n$ and not $2n$]
If the array shrank as soon as $\cc{cap}\ge2n$, then right after a doubling
($\cc{cap}=2n$) one removal would trigger a shrink and one insertion a grow: every
operation would cost $\Theta(n)$. Waiting until $\cc{cap}\ge3n$ leaves a gap of
$\Omega(n)$ operations between consecutive resizes, which restores $O(1)$ amortised
cost. Growing by a constant ($\cc{cap}+10$) instead of doubling copies
$10+20+30+\dots$ items, $\Theta(m^2)$ in total.
\end{remark}

\begin{figure}[ht]\centering
\begin{tikzpicture}[x=7.5mm]
\foreach \v [count=\i from 0] in {3,1,4,1} \node[cell] (o\i) at (\i,0) {\v};
\node[left,font=\scriptsize,align=right] at (-0.6,0) {size 4,\\ cap.\ 4};
\foreach \v [count=\i from 0] in {3,1,4,1,5,{},{},{}} \node[cell] (n\i) at (\i,-1.5) {\v};
\node[cell,sh] at (4,-1.5) {5};
\node[left,font=\scriptsize,align=right] at (-0.6,-1.5) {size 5,\\ cap.\ 8};
\foreach \i in {0,...,3} \draw[ptr,dashed] (o\i.south) -- (n\i.north);
\node[right,font=\scriptsize] at (4.6,0) {old array freed};
\node[right,font=\scriptsize] at (4.6,-0.75) {copy: $\Theta(n)$};
\node[right,font=\scriptsize] at (8.2,-1.5) {$\leftarrow$ \texttt{push\_back(5)}};
\end{tikzpicture}
\caption{\cc{push\_back} on a full vector: a new array of double capacity is allocated,
the items are copied, the old array is freed, then 5 is stored. All pointers and
iterators into the old array are invalid afterwards.}
\end{figure}

\section{Sorted array lists and binary search}\label{sec:bsearch}

If the list is kept in non-decreasing order, searching need not be linear.

\begin{definition}[Lower and upper bound]
$\mathrm{lb}(x)$ is the smallest index $i$ with $\cc{items}[i]\ge x$, or $n$ if none;
$\mathrm{ub}(x)$ is the smallest index $i$ with $\cc{items}[i]>x$, or $n$ if none.
\end{definition}

Then $x$ occurs in the list iff $\mathrm{lb}(x)<n$ and $\cc{items}[\mathrm{lb}(x)]=x$;
the number of occurrences is $\mathrm{ub}(x)-\mathrm{lb}(x)$; and the sorted position for
inserting $x$ \emph{after} equal items is $\mathrm{ub}(x)$.

\listingfile{listings/lowerbound.c}{Lower bound by bisection.}

\begin{theorem}[Correctness and cost of bisection]
The invariant holds initially (the ranges are empty), is preserved by both branches, and
when $\mathit{lo}=\mathit{hi}$ it says exactly that $\mathit{lo}=\mathrm{lb}(x)$. Each
round halves $\mathit{hi}-\mathit{lo}$, so the loop runs at most $\lfloor\lg n\rfloor+1$
times: search is $\Theta(\log n)$.
\end{theorem}

\begin{figure}[ht]\centering
\begin{tikzpicture}[x=7.5mm,y=9mm]
\foreach \row/\lo/\hi/\mid/\txt in {0/0/8/4/{$lo=0,\ hi=8,\ mid=4$: \ $a[4]=16<23$, so $lo=5$},
  1/5/8/6/{$lo=5,\ hi=8,\ mid=6$: \ $a[6]=38\ge23$, so $hi=6$},
  2/5/6/5/{$lo=5,\ hi=6,\ mid=5$: \ $a[5]=23\ge23$, so $hi=5$}} {
  \foreach \v [count=\i from 0] in {2,5,8,12,16,23,38,56} {
    \pgfmathtruncatemacro{\out}{(\i<\lo)||(\i>=\hi)}
    \ifnum\out=1 \node[cell,sh] at (\i,-\row) {\v}; \else \node[cell] at (\i,-\row) {\v}; \fi }
  \node[cell,shd] at (\mid,-\row) {\pgfmathparse{{2,5,8,12,16,23,38,56}[\mid]}\pgfmathresult};
  \node[right,font=\scriptsize] at (8.1,-\row) {\txt};
}
\foreach \i in {0,...,7} \node[idx] at (\i,0.48) {\i};
\node[font=\scriptsize] at (3.5,-2.9) {$lo=hi=5$: the lower bound of 23 is index 5};
\end{tikzpicture}
\caption{\cc{lowerBound(a, 8, 23)}. Shaded cells are already excluded; the dark cell is
\cc{mid}. The range $[lo,hi)$ halves each round.}
\end{figure}

\begin{proposition}[Cost of the sorted array list]
\cc{search}, \cc{count}, \cc{lowerBound}, \cc{upperBound}: $\Theta(\log n)$. \cc{insert}
and \cc{erase}: $\Theta(\log n)$ to locate plus $\Theta(n)$ to shift, i.e.\ $\Theta(n)$
overall. Sorting does not make \emph{insertion} cheaper; it makes \emph{searching}
cheaper. Complete code appears in Problem~\ref{prob:4}.
\end{proposition}

\section{Index techniques for array problems}\label{sec:arraytech}

Most array problems of the online-judge kind are solved by one of four patterns. Each
replaces a quadratic double loop by a single pass.

\para{Read index and write index} One index $r$ reads every item; a second index $w$ marks
the end of the part already built. An item that must be kept is copied to \cc{a[w++]}.
Because $w\le r$ at all times, nothing unread is overwritten. Used for removing
duplicates, moving zeroes, in-place compression (Problems~\ref{prob:1},~\ref{prob:7},~\ref{prob:10}).

\para{Two indices from the two ends} On a \emph{sorted} array, \cc{lo} starts at the left
and \cc{hi} at the right; at each step one of them moves inward, and the sortedness proves
that the discarded index cannot be part of an answer. Used for two-sum and for palindrome
checks (Problem~\ref{prob:8}; on a doubly linked list, Problem~\ref{prob:20}).

\para{The reversal trick} Reversing a block is a two-ends loop. Rotation by $k$ is three
reversals, transpose followed by row reversal rotates a matrix
(Problems~\ref{prob:5},~\ref{prob:12}).

\para{Running summaries} One left-to-right pass keeps a summary of the prefix read so
far: a running sum, the best sum of a subarray ending here (Kadane, Problem~\ref{prob:9}),
or a vote counter (Boyer--Moore, Problem~\ref{prob:11}). Chapters~\ref{ch:stacks}
and~\ref{ch:deques} extend this idea: the summary becomes a \emph{stack} or a
\emph{deque} of candidate indices.
